
This work investigated whether dataset metadata completeness predicts reproducibility variance before experiments run. Using synthetic data modeled on OpenML benchmark datasets with matched experimental configurations, results indicate that top-quartile metadata completeness is associated with a 42.1\% reduction in performance variance (IQR), with preprocessing entropy accounting for 64.7\% of this association.

The proposed mechanism is epistemic entropy reduction: complete documentation constrains the degrees of freedom available to implementing researchers. When metadata specifies preprocessing steps, missing value handling, and train/test splits, researchers may converge on similar pipelines, producing more consistent results.

Limitations warrant acknowledgment. All validation used synthetic data due to API availability constraints; real-world replication is necessary before drawing conclusions about actual datasets. The design is observational---the work predicts but cannot establish causation. Results are specific to synthetic data modeled on OpenML tabular datasets; generalization requires separate validation.

The broader implication, if findings replicate with real data, would be a reframing of reproducibility---not as a binary property of individual papers to be assessed post-hoc, but as a continuous property of datasets to be predicted proactively.

